\section{Conclusion} %
\label{sec:conclusion}

This paper presented context types, a new kind of type refinement in
which type qualifiers are equipped with a modality indicating whether
it should be interpreted as an over- or under-approximate
specification. The resulting type system permits both kinds of
specifications to be compositionally combined within a unified
framework for safety and reachability verification. The key insight is
to treat a typing context as a contextual capability --- a resource
whose structure governs the demonic and angelic choices available when
reasoning about terms. Building on this insight, we have developed a
resource algebra of capabilities and a modal context logic over that
algebra, and given a denotation of our type system into context logic.
Bunched context combinators track the dependencies between variables
and determine when angelic choices may be made independently, the
context sum operator recovers safety guarantees along executions that
reachability assumptions ensure exist, and a persistence modality
extends these guarantees to nondeterministic programs. Our metatheory,
mechanized in Rocq, establishes that a single soundness theorem
delivers both guarantees: every value produced by a well-typed term
satisfies its demonic types, and every value prescribed by its angelic
types is witnessed by some execution. Case studies drawn from
compositional may/must analysis, coverage verification, and
incorrectness reasoning suggest that this unified foundation can
tackle verification tasks previously addressed by separate,
property-specific systems.
